\section{Introduction}
\label{sec:introduction}

Drought frequency and severity have increased worldwide, driven by
ongoing climate change and rising anthropogenic water demands
\citep{IPCC2023, Pokhrel2021}, placing growing pressure on surface water
storage systems. In the Mediterranean basin, these trends are particularly
acute: multi-year dry spells have repeatedly pushed reservoir stocks to
critically low levels, threatening agricultural supply, hydropower
generation, and urban water distribution \citep{Cammalleri2022,
Bevacqua2024}. Sicily offers one of the most illustrative expressions of
this vulnerability: over the past decade, its reservoirs have repeatedly
reached critically low levels during the dry season \citep{ISPRA2023},
and total stored volumes have rarely recovered to historical averages even
after rainfall events, signalling a progressive depletion trend in
regional water reserves.

Accurate and continuous monitoring of reservoir storage is therefore
essential for adaptive water management. The principal sources of
information are in-situ gauge networks and remote-sensing time series,
each with specific strengths and limitations. Water-level gauges provide
near-real-time data, but converting levels into volumes requires
area--elevation--volume (AEV) relationships. In many systems, these still
rely on design-phase bathymetric surveys (1950s--1960s), which may not
reflect current conditions due to sedimentation and morphological change
\citep{BrunoEtAl2003}. Satellite radar altimetry --- as provided, for
example, by the DAHITI database \citep{Schwatke2015} --- offers a
complementary and globally consistent source of water-level time series,
now covering several thousand lakes and reservoirs worldwide. Combined
with global reservoir attribute databases such as GRanD
\citep{Lehner2011} and GRDL \citep{Hao2024}, which compile AEV
relationships for over 7,000 water bodies, and in-situ and altimetry-based
records compiled in GROWL \citep{Zhang2026}, these platforms have
substantially expanded the observational basis for large-scale hydrological
assessments. It is important to note, however, that converting any
remotely sensed water-level or surface-area signal into a storage estimate
still requires AEV information --- whether derived from in-situ surveys,
bathymetric modelling, or the design-phase documentation that many
reservoirs rely on. The present study is no exception: the storage time
series presented here are also derived from historical AEV relationships,
and the associated uncertainty is explicitly acknowledged as a limitation.

Alongside altimetry, Synthetic Aperture Radar (SAR) imagery offers a
complementary approach to reservoir monitoring based on the direct
delineation of the water surface extent. Unlike optical sensors operating
in the visible and near-infrared range, SAR is not affected by cloud cover
or solar illumination \citep{Torres2012, Mullissa2021}, a critical
advantage in Mediterranean and tropical contexts where cloud contamination
can severely reduce the temporal availability of optical observations
\citep{Pekel2016}. Processed on cloud computing platforms such as Google
Earth Engine (GEE), Sentinel-1 SAR data enable systematic surface water
mapping over large areas and long time periods \citep{Gorelick2017}.
Recent studies confirm this operational maturity: \citet{BhatpuriaEtAl2024}
applied automated SAR workflows for storage monitoring in semi-arid India;
\citet{ValerioEtAl2024} showed that Sentinel-1 combined with Sentinel-2 on
GEE reliably detects water bodies smaller than 0.5~ha; and
\citet{PengEtAl2025} demonstrated GEE-based SAR--optical fusion for flood
mapping with accuracies above 92\%. Relative to altimetry, area-based SAR
approaches have the advantage of providing spatially explicit water masks
that can capture asymmetric drawdown patterns, detect disconnected water
pools, and function independently of whether a satellite altimetry track
passes over the water body. The two approaches are therefore
complementary: altimetry constrains water levels with centimetric accuracy
where coverage exists, while SAR provides spatial and temporal coverage
of the water surface in systems not reached by altimetry tracks.

Despite this progress, SAR-based reservoir delineation exhibits
substantial variability in performance across water bodies that cannot be
attributed to sensor characteristics alone. Wind-driven surface roughness
raises backscatter over open water, causing misclassification as non-water
\citep{Bioresita2019, CazalsEtAl2021}. At the water--land boundary,
mixed pixels introduce systematic errors whose severity depends on the
ratio of shoreline length to water body area \citep{CazalsEtAl2021}.
Shadowing and layover effects in topographically complex settings further
complicate delineation near steep reservoir margins \citep{Martinis2015}.
These phenomena are well-established individually; yet a key operational
question remains unanswered: \emph{which geometric property of a
reservoir most efficiently predicts its SAR classification skill, and can
it be estimated before any SAR processing is performed?}

We propose that the answer lies in the \textbf{area-to-perimeter ratio
(A/P; units: m)}, a length-scale index that captures the compactness of
the shoreline relative to the water body's spatial extent. The physical
reasoning proceeds as follows. For a sensor with ground resolution $r$,
the boundary strip of width $r$ enclosing the perimeter $P$ contains
approximately $P \cdot r$ square metres of ambiguous pixels (neither
purely water nor purely land). As a fraction of total water area $A$,
this amounts to $P \cdot r / A = r / (A/P)$, a quantity that grows as
$A/P$ decreases --- i.e.\ as the shoreline becomes more irregular or
elongated relative to the enclosed area. This geometric argument is
corroborated by \citet{CazalsEtAl2021}, who quantified a systematic
perimeter-inward erosion of $\sim$20~m in Sentinel-1-derived water masks
at the pixel level. Beyond the mixed-pixel mechanism, elongated reservoirs
(low A/P) also present longer wind-fetch distances along their main axis,
which promotes wave development and elevated backscatter that can
trigger omission errors \citep{McQuillan2025}. Compact reservoirs (high
A/P), by contrast, expose shorter fetch lengths, simpler shoreline
geometries, and fewer radar shadow or layover zones. Together, these
mechanisms predict that A/P is a reliable a-priori indicator of SAR
classification performance --- one that, to our knowledge, has not been
systematically tested across a multi-reservoir dataset.

This paper addresses that gap. Using a GEE-based workflow combining
Sentinel-1 imagery and Support Vector Machine (SVM) classification across
41 reservoirs in Sicily, we: (i)~quantify the statistical relationship
between A/P and SAR classification accuracy; (ii)~validate the resulting
surface-area and storage time series against high-resolution PlanetScope
imagery and official storage records; and (iii)~demonstrate that, when
training samples are drawn automatically from the JRC Global Surface Water
dataset \citep{Pekel2016}, the same workflow can be transferred to
reservoirs in diverse climatic regions worldwide, with performance within
0.05--0.07 KGE of the manually supervised baseline. The A/P ratio emerges
from this analysis as a practical screening tool: computable from global
databases such as GRanD or GROWL before any SAR processing, it allows
practitioners to anticipate where the workflow will perform well and where
additional caution or complementary altimetry data are warranted.

The remainder of the paper is organised as follows.
Section~\ref{sec:study_area} describes the study area and the 41 Sicilian
reservoirs. Section~\ref{sec:methods} details the SAR preprocessing, SVM
classification, morphometric analysis, and validation strategy.
Section~\ref{sec:results} presents the classification results, the
A/P--performance relationship, and the global transferability assessment.
Section~\ref{sec:discussion} discusses implications for operational
monitoring, the limitations of the approach --- including the dependence
on historical AEV relationships --- and avenues for future work.
Section~\ref{sec:conclusions} summarises the main conclusions.
